% ============================================================
%  Section 1 -- Introduction
% ============================================================

\chapter{Introduction}

\section{Contexte industriel}

Les assistants fondes sur les grands modeles de langue sont devenus des
candidats naturels pour automatiser une partie du support informatique :
diagnostic initial, recherche de procedure, reformulation de consignes ou
orientation vers le bon niveau d'escalade. Dans un contexte industriel comme
Michelin, cette promesse doit cependant etre confrontee a des contraintes
strictes : fiabilite des reponses, tracabilite, protection des donnees internes,
cout d'execution, latence acceptable et impact environnemental.

Le risque principal n'est pas seulement qu'un modele reponde mal. Il est qu'une
reponse plausible, non verifiee, soit integree dans un processus operationnel.
Un ticket IT contient souvent des indices techniques incomplets, parfois des
donnees sensibles, et peut necessiter une decision fondee sur un etat externe :
statut d'un service, procedure interne, incident connu, documentation ou niveau
d'urgence. Une architecture purement conversationnelle peut produire une reponse
fluide mais deconnectee de ces sources.

\bibopsname{} a ete concu pour etudier cette question sous forme experimentale.
Le projet ne cherche pas uniquement a construire un agent, mais a construire un
\emph{banc d'evaluation} capable de comparer plusieurs architectures et de
produire des artefacts exploitables : fichiers JSON, graphiques, traces
d'execution, scores composites et rapports de securite. Cette orientation MLOps
est centrale : l'objectif est de rendre les decisions mesurables plutot que de
se fier a une impression qualitative.

\section{Objectifs du projet}

Le projet fixe trois objectifs principaux.

\begin{enumerate}[label=\textbf{O\arabic*.}]
  \item \textbf{Comparer deux architectures de support IT.} La premiere est un
  LLM unique, interroge sans outil. La seconde est un systeme multi-agents
  simplifie, centre sur une boucle ReAct et trois outils : verification de
  serveur, recherche dans une base de connaissances et recherche documentaire
  vectorielle.

  \item \textbf{Construire une evaluation multi-criteres.} Les sorties ne sont
  pas jugees seulement sur leur qualite. Le framework agregue qualite,
  securite, cout, latence et empreinte environnementale dans une politique
  composite explicite.

  \item \textbf{Tester les risques agentiques.} La Racing Arena simule un
  environnement distribue dans lequel plusieurs agents recoivent des flux
  temps reel et peuvent etre attaques. Elle sert a observer les effets de
  l'outillage, des APIs faibles et des injections de prompt dans un contexte
  dynamique.
\end{enumerate}

Ces objectifs expliquent la structure du depot. \file{src/agent/} contient
l'agent IT, \file{src/bibops/evaluation/} contient le moteur d'evaluation,
\file{src/benchmark/} et \file{scripts/benchmark/} orchestrent les campagnes, et
\file{src/racing/} porte l'arene distribuee.

\section{Perimetre du rapport}

Le present rapport analyse le depot tel qu'il existe le 12 mai 2026. Les
chiffres sont extraits des artefacts locaux deja produits :
\file{comparison_results.json} pour la comparaison principale,
\file{security_race_report.json} pour la Racing Arena,
\file{benchmark_mcp.json} et \file{benchmark_mcp_tools.json} pour les outils
MCP, ainsi que les rapports A/B, position bias, A2A et Kaggle.

Le rapport suit le plan impose par l'ecole :

\begin{itemize}
  \item description conceptuelle du probleme ;
  \item presentation de la solution ;
  \item resultats experimentaux avec figures ;
  \item discussion des limites ;
  \item conclusion et perspectives.
\end{itemize}

La demarche reste volontairement factuelle. Lorsque les donnees disponibles ne
permettent pas de conclure, le rapport le signale explicitement. C'est le cas,
par exemple, du benchmark A2A du 12 mai 2026 : le pipeline est present et les
artefacts existent, mais les neuf agents testes remontent chacun une erreur et
ne permettent donc pas de tirer une conclusion de performance.

\section{Apports principaux}

Les apports de \bibopsname{} sont les suivants.

\begin{itemize}
  \item Un agent ReAct testable, dont les appels outils sont bornes par des
  politiques de timeout et de retry.
  \item Une separation claire entre production, recherche, benchmarks et tests.
  \item Une evaluation composite explicite, avec poids et seuils visibles dans
  le code.
  \item Une instrumentation des traces, permettant de relier une reponse finale
  a ses appels outils et a son niveau de confiance.
  \item Une arene de securite agentique qui rend observables les injections, les
  fuites et les comportements d'agents en temps reel.
\end{itemize}

Le resultat le plus important est nuance : l'architecture multi-agents est
meilleure que le LLM unique sur la campagne principale, mais elle n'atteint pas
encore le seuil de qualite absolu fixe par la politique de release. Le projet
montre donc a la fois l'interet de l'approche outillee et la necessite d'un
travail supplementaire avant industrialisation.
